\section{High Score}
\label{sec:cses-high-score}

\problemstatement{A directed graph of $n$ rooms and $m$ edges, each edge
carrying a score (possibly negative). Starting in room $1$ and ending in
room $n$, maximise the total collected score. If the score can be made
arbitrarily large, output $-1$.}

\begin{patternbox}
Maximising with possibly negative (and positive) edges, plus detecting an
unbounded answer, is \textbf{Bellman--Ford} run for maxima. A positive
cycle that lies on some $1\to n$ path makes the score unbounded -- the
analogue of a negative cycle in the minimisation form.
\end{patternbox}

\subsection*{Maximise, then check reachable positive cycles}

Set $\textit{dist}[1]=0$, others $-\infty$, and relax every edge $n-1$
times using $\max$: $\textit{dist}[v]=\max(\textit{dist}[v],
\textit{dist}[u]+w)$. After that, run one more round: any node whose value
still increases sits on or downstream of a positive cycle. Mark such nodes
and propagate the mark forward; if node $n$ ends up marked \emph{and} is
reachable from $1$, the answer is $-1$. Otherwise $\textit{dist}[n]$ is the
maximum score.

\begin{ideabox}[A Positive Cycle on the Path Means Unbounded]
Only cycles that are both reachable from $1$ and can reach $n$ matter --
looping such a cycle adds score forever. Extra relaxation rounds detect the
still-improving nodes, and a forward propagation checks whether $n$ is
affected.
\end{ideabox}

\subsection*{Algorithm}

\begin{enumerate}
  \item Relax all edges $n-1$ times with $\max$ (from $\textit{dist}[1]=0$).
  \item Run $n$ more rounds; if an edge still improves $v$, mark $v=+\infty$
        and propagate the mark along edges.
  \item If $\textit{dist}[n]=+\infty$ output $-1$, else output
        $\textit{dist}[n]$.
\end{enumerate}

\subsection*{Dry run}

\dryrun{$1\!\to\!2(5)$, $2\!\to\!3(-2)$, $3\!\to\!2(3)$ (cycle $2\!\to\!3
\!\to\!2$ gains $+1$), $3\!\to\!4(1)$; $n=4$.}

\begin{itemize}
  \item Cycle $2\!\to\!3\!\to\!2$ has net $+1$ and reaches $4$ -- score is
        unbounded.
  \item Output $-1$.
\end{itemize}

\subsection*{C++ implementation}

\begin{lstlisting}
#include <bits/stdc++.h>
using namespace std;
const long long NEG = LLONG_MIN / 4;

int main() {
    int n, m;
    cin >> n >> m;
    vector<array<long long,3>> edges(m);          // {u, v, w}
    for (auto& e : edges) cin >> e[0] >> e[1] >> e[2];

    vector<long long> dist(n + 1, NEG);
    dist[1] = 0;
    for (int i = 0; i < n - 1; i++)
        for (auto& [u, v, w] : edges)
            if (dist[u] > NEG) dist[v] = max(dist[v], dist[u] + w);

    // extra rounds: mark nodes on/after a positive cycle as +infinity
    for (int i = 0; i < n; i++)
        for (auto& [u, v, w] : edges)
            if (dist[u] > NEG && dist[u] + w > dist[v]) dist[v] = LLONG_MAX / 4;

    cout << (dist[n] >= LLONG_MAX / 8 ? -1 : dist[n]) << "\n";
    return 0;
}
\end{lstlisting}

\begin{complexitybox}
\textbf{Time Complexity.} $O(nm)$ -- Bellman--Ford's relaxation rounds.
\textbf{Space Complexity.} $O(n+m)$.
\end{complexitybox}

\begin{takeawaybox}
Maximising a path score with mixed-sign edges is Bellman--Ford with $\max$;
a reachable positive cycle makes it unbounded. Running extra rounds and
propagating a ``$+\infty$'' mark cleanly detects the nodes an unbounded
cycle can inflate -- the mirror of negative-cycle detection.
\end{takeawaybox}
